\section{Exact data}\label{app:data}

Every design recorded here is given by data from which
\eqref{eq:parseval} can be checked in finite arithmetic.  Radicals occur
in three places: the
icosahedral design of \S\ref{app:icosa}, the whitener of a graphic design
in \S\ref{app:diamond}, and the unit tight eigenvectors of the
tetrahedron.  Each time, the algebraic relation that eliminates them is
displayed.

\subsection{The tetrahedron at \texorpdfstring{$(4,3)$}{(4,3)}}\label{app:tetra}

Atoms and weights:
\[
  g_0=(1,1,1),\quad g_1=(1,-1,-1),\quad g_2=(-1,1,-1),\quad g_3=(-1,-1,1),
  \qquad t_c=\tfrac14 .
\]
These are the four sign vectors with an even number of minus signs.  For
$i\ne j$ the coordinatewise products of $g_i$ and $g_j$ sum to $-1$, and for
each pair of distinct coordinates the four products $(g_c)_a(g_c)_b$ are
two $+1$ and two $-1$; hence
\[
  \ell_c=3,\qquad \langle g_i,g_j\rangle=-1 \ (i\ne j),
  \qquad \sum_c g_c^{\vphantom{\mathsf{T}}}g_c\tp=4I_3 ,
\]
so $\sum_c t_c g_c^{\vphantom{\mathsf{T}}}g_c\tp=I_3$ and this is a design.  It is the
design of Example~\ref{ex:tetra}, atom for atom and weight for weight.
Lemma~\ref{lem:trace} holds termwise: $t_c\ell_c=\tfrac34$ and
$\sum_c t_c\ell_c=3=k$.

For $C=\{0,1,2,3\}\setminus\{j\}$,
\[
  \SC=4I_3-g_j^{\vphantom{\mathsf{T}}}g_j\tp,
  \qquad
  \operatorname{spec}\SC=\{1,4,4\},
  \qquad
  \operatorname{spec}(\SC-I)=\{0,3,3\},
\]
the eigenvalue $1$ occurring along $g_j$.  For $j=3$ the gap is the singular
positive form
\[
  x\tp(S_{\{0,1,2\}}-I)x=(x_0-x_1)^2+(x_0+x_2)^2+(x_1+x_2)^2 ,
\]
with null line $\R\,(1,1,-1)$.  Every triple dominates and none dominates
strictly: the tetrahedron is extremal in the sense of
Definition~\ref{def:tie}.

\emph{Chart and volume sampling.}  Here $v_c=\tfrac12 g_c$, so
$P_{cd}=\tfrac14\langle g_c,g_d\rangle$, that is $P=\tfrac14(4I_4-J_4)$ with
$J_n$ the all-ones $n\times n$ matrix.  For a triple $C$,
$P[C]=\tfrac14(4I_3-J_3)$, whose eigenvalues are $\tfrac14\{1,4,4\}$, so by
\eqref{eq:dpp}
\[
  \pi_C=\det P[C]=\tfrac{16}{64}=\tfrac14 \qquad\text{for each of the four triples,}
\]
confirming $\sum_C\pi_C=1$.  All four characteristic polynomials coincide,
whence
\[
  q(x)=\E_\pi[\chi_{\SC}(x)]=(x-1)(x-4)^2=x^3-9x^2+24x-16 ,
\]
\[
  q(1)=0,\qquad q'(1)=9,\qquad q''(1)=-12 .
\]
The subleading coefficient is $-9=-\sum_c t_c\ell_c^2=-4\cdot\tfrac14\cdot9$,
attaining the bound $-k^2$ of Proposition~\ref{prop:qone}.

\emph{Stationarity data.}  Take $\lambda_C=\tfrac14$ on each of the four
triples and $u_C=g_{c_0}/\sqrt3$ for $C$ the complement of $\{c_0\}$.  Then
\[
  \Lambda=\cv\sum_C\lambda_C\,u_C^{\vphantom{\mathsf{T}}}u_C\tp
   =1\cdot\sum_{c_0}\tfrac14\cdot\tfrac13\,g_{c_0}^{\vphantom{\mathsf{T}}}g_{c_0}\tp
   =\tfrac1{12}\cdot4I_3=\tfrac13 I_3,
  \qquad \tr\Lambda=1=\cv .
\]
The quadric law \eqref{eq:quadric} reads $g_c\tp\Lambda g_c=\ell_c/3=1=\cv$, and
the coverage law \eqref{eq:coverage} reads, for each $c$, a sum over the three
triples containing $c$ of $\tfrac14\bigl(-\tfrac1{\sqrt3}\bigr)^2=\tfrac1{12}$,
that is $\tfrac14=t_c\cv$.

\emph{Two moments.}  The Gram matrix of a triple is $4I_3-J_3$, of trace $9$
and determinant $16$, so the quantity $\gamma$ of \eqref{eq:twomoment} is
$4\cdot16-(9-1)^2=0$: the two-moment criterion of
Proposition~\ref{prop:twomoment} is tight here.

\subsection{The corank-one extremal family}\label{app:corankone}

Let $m=k+1$ and let $t$ be any point of the open simplex.  By
Theorem~\ref{thm:corankone} the extremal designs at $(k+1,k)$ with weights $t$
are exactly those whose chart is $P=I-zz\adj$ with
\[
  \lvert z_c\rvert^2=\frac{1-t_c}{k},
  \qquad\text{equivalently}\qquad
  s_c=P_{cc}=\frac{k-1+t_c}{k},
  \qquad
  \ell_c=\frac{k-1+t_c}{k\,t_c} .
\]
The normalization is automatic: $\sum_c(1-t_c)=m-1=k$.  Over $\R$ the vector
$z$ is determined up to signs and an orthogonal change of the ambient frame;
over $\C$, up to phases.  At $k=3$, $t\equiv\tfrac14$ one gets
$\lvert z_c\rvert^2=\tfrac14$, $s_c=\tfrac34$, $\ell_c=3$ --- the tetrahedron of
\S\ref{app:tetra}.  At $k=3$, $t=(\tfrac23,\tfrac19,\tfrac19,\tfrac19)$ one
gets a member which is not a rotation of the tetrahedron:
\[
  g_0=\tfrac13(2,2,2),\quad
  g_1=\tfrac13(2,2,-7),\quad
  g_2=\tfrac13(-7,2,2),\quad
  g_3=\tfrac13(2,-7,2),
\]
with $\ell=(\tfrac43,\tfrac{19}3,\tfrac{19}3,\tfrac{19}3)$, the single
dependency $3g_0+2g_1+2g_2+2g_3=0$, and dominating triple $\{1,2,3\}$ whose
gap is $\tfrac83\bigl[(x_0-x_1)^2+(x_0-x_2)^2+(x_1-x_2)^2\bigr]$.  The
equality locus is therefore not a single orbit even at corank one.  Two
readings of this member.  Its squared cosines take two values,
$\tfrac1{19}$ between $g_0$ and each other atom and $\tfrac{64}{361}$
between the others, so a real extremal design need not be equiangular
(\texttt{Gtz.sharp\allowbreak Atom\_not\_equiangular}) --- against the icosahedron of
\S\ref{app:icosa}, whose single squared cosine is $\tfrac15$.  And the
dependency coefficients $y=(3,2,2,2)$ satisfy $y_c^2/\bigl(t_c(1-t_c)\bigr)
=\tfrac{81}2$ at every atom, which is \eqref{eq:tielaw} read on this
design.

The family has a graphic realization, uniform in the rank: the
$(k+1)$-cycle at unit conductance and uniform weight, whose dominating
subsets are its paths, dominating with margin exactly zero
(\texttt{Gtz.cycle\allowbreak Graph\allowbreak Data\_dominates},
\texttt{Gtz.cycle\allowbreak Graph\allowbreak Data\_not\_pos\allowbreak Def}).  The gap is a telescoping
Cauchy--Schwarz, so the vanishing of the margin is visible at every rank
rather than instance by instance.

\subsection{The chart point of
\texorpdfstring{\S\ref{sec:sharp}}{the sharpness section}}\label{app:liftfail}

Fix the symmetric matrix
\begin{equation}\label{eq:liftchart}
  P=\frac13\begin{pmatrix}
    2&1&0&1\\ 1&1&1&0\\ 0&1&2&-1\\ 1&0&-1&1\end{pmatrix},
  \qquad M:=3P .
\end{equation}
Then $M^2=3M$, as one checks entry by entry, hence $P^2=P$, and
$\tr P=\tfrac{2+1+2+1}{3}=2$.  So
$P$ is a symmetric idempotent of rank two and $(P,t)$ is a chart point of size
four and rank two in the sense of Definition~\ref{def:chartpoint} for every $t$
in the closed simplex.  Take the escaping index $z=1$ and the one-parameter
family of weights
\begin{equation}\label{eq:liftweights}
  t^{(\tau)}=\Bigl(\tfrac{1-\tau}{3},\ \tau,\ \tfrac{1-\tau}{3},\
  \tfrac{1-\tau}{3}\Bigr),\qquad 0<\tau<1 ,
\end{equation}
whose member at $\tau=\tfrac14$ is the uniform vector.  With
$W=P-\diag(t^{(\tau)})$,
\[
  W[\{0,1\}]=
  \begin{pmatrix}\dfrac{1+\tau}{3} & \dfrac13\\[6pt]
                 \dfrac13 & \dfrac13-\tau\end{pmatrix},
  \qquad
  \det W[\{0,1\}]
   =\frac{(1+\tau)(1-3\tau)-1}{9}
   =-\frac{\tau(2+3\tau)}{9} .
\]
The determinant is strictly negative for every $\tau\in(0,1)$ and vanishes to
first order as $\tau\to0$.  An explicit witness is $(-1,\,1+\tau)$:
\[
  (-1,1+\tau)\,W[\{0,1\}]\,(-1,1+\tau)\tp
  =\frac{1+\tau}{3}\bigl[1-2+(1+\tau)(1-3\tau)\bigr]
  =-\tfrac13(1+\tau)\,\tau\,(2+3\tau).
\]
At $\tau=\tfrac14$ the block, its determinant and the probe $(1,-4)$ are
\[
  W[\{0,1\}]=\begin{pmatrix}5/12&1/3\\ 1/3&1/12\end{pmatrix},
  \qquad \det=-\tfrac{11}{144},
  \qquad (1,-4)\,W[\{0,1\}]\,(1,-4)\tp=-\tfrac{11}{12}.
\]
The leverages $P_{cc}/t_c$ at $\tau=\tfrac14$ are
$\tfrac83,\tfrac43,\tfrac83,\tfrac43$, all exceeding $1$, so the point is not
disposed of by deletion of a light atom.

The deflation \eqref{eq:deflate} at $z=1$, where $P_{11}=\tfrac13>0$, is a
symmetric idempotent of trace $1$ with $P'_{00}=\tfrac13$.  Deleting $z$ and
renormalizing puts each survivor at weight
$\bigl((1-\tau)/3\bigr)/(1-\tau)=\tfrac13$, so the compressed gap at the
surviving index $0$ is exactly $0$, at every $\tau$.

Theorem~\ref{thm:sharp} exploits that frozen equality: the Schur term of
\eqref{eq:schur} is matched in full at $t_z=0$ and falls short by a quantity
of first order in $t_z$ thereafter, with no downstairs slack to close the
gap.

\begin{wip}
Exhibit a one-parameter deformation of \eqref{eq:liftchart} with compressed
slack $\varepsilon>0$, and display the resulting determinant
$(3\varepsilon-2\tau-3\tau^{2})/9$ of \eqref{eq:family} changing sign at
$\tau\asymp \varepsilon$.
\end{wip}

\subsection{Graphic designs, and the diamond at
\texorpdfstring{$(5,3)$}{(5,3)}}\label{app:diamond}

The diamond is the second extremal design at rank three.  Let $\gph$ be a
connected multigraph on a vertex set $V$ with a distinguished vertex, the
\emph{ground}; let $b_e\in\R^{V\setminus\{\mathrm{ground}\}}$ be the reduced
incidence vector of the edge $e$ ($+1$ at its tail, $-1$ at its head, the
ground coordinate deleted); let $c_e>0$ be conductances and $t_e>0$ weights
summing to one.  Put
\[
  L:=\sum_e c_e\,b_e^{\vphantom{\mathsf{T}}}b_e\tp,
\]
the weighted reduced Laplacian, positive definite because $\gph$ is connected.
Choose $R$ with $R\tp LR=I$ and set
\begin{equation}\label{eq:graphicatom}
  g_e:=\sqrt{c_e/t_e}\;R\tp b_e .
\end{equation}
Then $\sum_e t_e g_e^{\vphantom{\mathsf{T}}}g_e\tp=R\tp LR=I$, so this is a design of
size $\lvert E\rvert$ and rank $\lvert V\rvert-1$.  From $R\tp LR=I$ one reads
$RR\tp=L^{-1}$, whence
\begin{equation}\label{eq:graphicleverage}
  s_e=t_e\ell_e=c_e\,b_e\tp L^{-1}b_e=c_e\,R_{\mathrm{eff}}(e) :
\end{equation}
the leverage score of an edge is its conductance times its effective
resistance, and Lemma~\ref{lem:trace} becomes
$\sum_e c_eR_{\mathrm{eff}}(e)=\lvert V\rvert-1$.  Congruence by the invertible
$R$ removes the whitener from \eqref{eq:dominate}: an edge set $E'$ with
$\lvert E'\rvert=\lvert V\rvert-1$ dominates if and only if
\begin{equation}\label{eq:graphicdominate}
  \sum_{e\in E'}\frac{c_e}{t_e}\,b_e^{\vphantom{\mathsf{T}}}b_e\tp\;\psd\;L .
\end{equation}
The atoms need not be rational and are determined only up to the orthogonal
group; at rational conductances the graph, $L$, the effective resistances and
the leverages are rational.

Take $\gph=K_4-e$ on vertices $a,b,c,d$ with $d$ the ground and $cd$ the missing
edge.  The five edges and their reduced incidence rows in the coordinates
$(y_a,y_b,y_c)$ are
\[
  \begin{array}{llllll}
    e_0=ab: & (1,-1,0), &\quad e_1=ac: & (1,0,-1), &\quad e_2=ad: & (1,0,0),\\
    e_3=bc: & (0,1,-1), &\quad e_4=bd: & (0,1,0). &&
  \end{array}
\]
Conductances $c=(2,3,3,3,3)$ and uniform weights $t_e=\tfrac15$.  Then
\[
  L=\begin{pmatrix}8&-2&-3\\ -2&8&-3\\ -3&-3&6\end{pmatrix},
  \qquad \det L=180,
  \qquad L^{-1}=\frac1{180}\begin{pmatrix}39&21&30\\ 21&39&30\\ 30&30&60
  \end{pmatrix},
\]
so $R_{\mathrm{eff}}(ab)=\tfrac1{180}(39-21-21+39)=\tfrac15$ and
$R_{\mathrm{eff}}(e)=\tfrac{39}{180}=\tfrac{13}{60}$ for each of the four rim
edges $ac,ad,bc,bd$.  By \eqref{eq:graphicleverage},
\[
  s_{ab}=\tfrac25,\quad s_{\mathrm{rim}}=\tfrac{13}{20},
  \qquad
  \ell_{ab}=2,\quad \ell_{\mathrm{rim}}=\tfrac{13}{4},
  \qquad
  \tfrac25+4\cdot\tfrac{13}{20}=3=k .
\]
Every leverage exceeds $1$.  The star at $a$, that is $E'=\{ab,ac,ad\}$,
dominates: the left side of \eqref{eq:graphicdominate} is
$5\bigl(2b_0b_0\tp+3b_1b_1\tp+3b_2b_2\tp\bigr)$, and subtracting $L$ gives
the form $Q$.  Doubling and completing the square in $y_a$ leaves $9y_c^2$:
\[
  2\,Q(y)\;=\;(-8y_a+2y_b+3y_c)^2+9y_c^2 ,
\]
positive semidefinite and singular, with null line $\R\,(1,4,0)$.  No
three-subset dominates strictly.  The eight spanning trees each carry a null
vector of their gap; the two circuits do worse.  At $E'=\{ab,ac,bc\}$ the
potential $y=(1,1,1)$ annihilates all three selected drops while
$y\tp Ly=3+3=6$, the two ground edges $ad$ and $bd$ contributing, so the gap
form equals $-6$; likewise at $E'=\{ab,ad,bd\}$ with $y=(0,0,1)$.  The
certificates for all ten subsets are listed in the proof of
Example~\ref{ex:diamond}.
The diamond is therefore an extremal design at $(5,3)$.  Its five incidence
rows are pairwise non-proportional, and by \eqref{eq:graphicatom}
proportionality of two atoms forces proportionality of the corresponding rows;
so no two atoms of the diamond are parallel.  This is
Example~\ref{ex:diamond}.

The conductances are forced: keep the rim conductance at
$c_{\mathrm{rim}}=3$, let the spine carry $c_{ab}=p>0$, and ask that the star
at $a$ dominate.  With $L$ recomputed at those conductances,
\eqref{eq:graphicdominate} for $E'=\{ab,ac,ad\}$ reads $Q\psd0$ for a
symmetric $3\times3$ form whose leading principal minors are $4p+24$,
$72(p-2)$ and $324p-648$.  For $p<2$ the second is negative, so $Q$ is not
positive semidefinite; for $p>2$ all three are positive, so $Q\pd0$; at
$p=2$ the form is the singular $Q$ above.  So $p=2$ is the unique conductance
at which the star dominates without dominating strictly, which is
extremality.  Then $K_4-e$ is three internally disjoint $a$-to-$b$ paths:
the edge $ab$ of conductance $2$ and two two-edge paths of conductance
$\tfrac32$ each, whose conductances add to $5$.  So
$R_{\mathrm{eff}}(ab)=\tfrac15$ and $s_{ab}=\tfrac25$, as computed
above.

Three facts about graphic designs in general, each formal.  Domination is a
matroid condition: an edge set of the right size dominates only if it is a
spanning tree, and positive definiteness of the Laplacian restricted to a
set is connectivity to the ground
(\texttt{Gtz.is\allowbreak Spanning\allowbreak Tree\_of\_dominates},
\texttt{Gtz.pos\allowbreak Def\_laplacian\allowbreak On\_iff\_is\allowbreak Ground\allowbreak Connected}), which is the
dictionary behind ``the eight spanning trees'' above.  Deleting a light edge
preserves domination in the Loewner order
(\texttt{Gtz.dominates\_of\_deleted\_dominates}), the graph reading of the
leverage ceiling.  And slack is available at the binding cell: $K_4$ at unit
conductance and uniform weight $\tfrac16$ is a design at $(6,3)$ whose star
at the ground dominates \emph{strictly}, with the explicit certificate
$2(y_a^2+y_b^2+y_c^2)+(y_a+y_b+y_c)^2$
(\texttt{Gtz.complete\allowbreak Four\allowbreak Data\_dominates\_strictly}).  It is the foil to the
diamond: same construction, one cell up, and the margin is positive.

The equality locus has a graphic realization at every rational point, and
with it a closed form for the leverage.  Bundle the $(k+1)$-cycle: replace
arc $j$ by $p_j$ parallel edges, and write $n=\sum_jp_j$.  Then an edge in
an arc of $p$ parallel edges has
\[
  \ell_e=\frac{(k-1)n+p}{k\,p},
  \qquad
  P_{ee}=\frac{k-1}{k\,p}+\frac{1}{k\,n}
  \qquad\text{(\texttt{Gtz.bundled\allowbreak Cycle\_leverage})},
\]
proved by exhibiting the transfer current, never by inverting the Laplacian,
and agreeing with \eqref{eq:tielaw} at the arc total $T=p/n$.  The gap of
the transversal omitting arc $j_0$ is
\[
  \sum_{j\ne j_0}\frac{d_j^2}{p_j}
  \;-\;\frac{\bigl(\sum_{j\ne j_0}d_j\bigr)^2}{\sum_{j\ne j_0}p_j}
  \qquad\text{(\texttt{Gtz.bundled\allowbreak Cycle\_gap\_eq\_engel\allowbreak Deficiency})},
\]
the Engel form of Cauchy--Schwarz in the drops $d_j$: nonnegative always,
and zero exactly on the ray $d_j\propto p_j$.  That is the mechanism behind
the extremality asserted in \S\ref{app:split} --- the margin vanishes
because a Cauchy--Schwarz is tight, not by an accident of the construction.
At rank two the gap is the exact rational square
$(d_1p_2-d_2p_1)^2/\bigl(p_1p_2(p_1+p_2)\bigr)$, so the slack off the tight
ray is an explicit function of the drops.

\subsection{Split families at
\texorpdfstring{$(6,3)$ and $(7,3)$}{(6,3) and (7,3)}}\label{app:split}

Splitting an atom into parallel copies preserves \eqref{eq:parseval}, because
the copies carry the same rank-one moment; and it preserves the value of the
objective, because a $k$-subset repeating a direction is singular while a
$k$-subset meeting $k$ distinct directions has the gap of the original.
Applied to the tetrahedron of \S\ref{app:tetra} this gives the two families
used throughout \S\ref{sec:ties} and Appendix~\ref{app:graveyard}.

\emph{At $(6,3)$.}  Directions $(0,1,2,2,3,3)$ on the four tetrahedron
vertices, weights
\[
  \bigl(\tfrac14,\ \tfrac14,\ a,\ \tfrac14-a,\ b,\ \tfrac14-b\bigr),
  \qquad 0<a,b<\tfrac14 ,
\]
so that every direction carries total weight $\tfrac14$.  Every member is
extremal.  Twelve triples attain $\lmin(\SC)=1$, those carrying three
distinct directions,
\[
  \begin{array}{cccccc}
   \{0,1,2\}&\{0,1,3\}&\{0,1,4\}&\{0,1,5\}&\{0,2,4\}&\{0,2,5\}\\
   \{0,3,4\}&\{0,3,5\}&\{1,2,4\}&\{1,2,5\}&\{1,3,4\}&\{1,3,5\}
  \end{array}
\]
missing the directions $3,3,2,2,1,1$ and $1,1,0,0,0,0$ respectively.

\emph{At $(7,3)$.}  Directions $(0,1,2,3,0,1,2)$, weights
$(\tfrac18,\tfrac18,\tfrac18,\tfrac14,\tfrac18,\tfrac18,\tfrac18)$; again each
direction carries $\tfrac14$, every leverage is exactly $3$, and
\eqref{eq:parseval} is exact.  Of the $35$ triples, the $20$ carrying three
distinct directions dominate, each with gap
$3I_3-g_d^{\vphantom{\mathsf{T}}}g_d\tp$ of spectrum $\{0,3,3\}$ where $d$ is the
missed direction; the remaining $15$ do not.  None of the $20$ has any slack.

More generally, fix a surjection $\mathrm{cl}$ from $\{1,\dots,m\}$ onto
$\{1,\dots,k+1\}$ and a weight vector $t$; give the index $c$ the corank-one
extremal atom of \S\ref{app:corankone} attached to the class
$\mathrm{cl}(c)$, evaluated at the class totals
$T_j:=\sum_{\mathrm{cl}(c)=j}t_c$.  Atoms in one class are equal, so the
fibrewise sums collapse and \eqref{eq:parseval} is inherited from the
corank-one design at $T$.  A $k$-subset repeating a class is singular; a
$k$-subset meeting $k$ classes has the gap of that corank-one design, which is
positive semidefinite and singular.  This is
Theorem~\ref{thm:tieseverywhere}.

At $k=3$ and $m=6$ the class profiles are $(3,1,1,1)$ and $(2,2,1,1)$; at
$m=7$ they are $(4,1,1,1)$, $(3,2,1,1)$ and $(2,2,2,1)$.  The two families
displayed above are the profiles $(2,2,1,1)$ and $(2,2,2,1)$.  Since a member
of this family has at most $k+1$ distinct directions and the diamond has
$k+2=5$, the diamond is not one of them.

\emph{A rational extremal design with an exact certificate.}  Take the $(6,3)$
split at $(a,b)=(\tfrac16,\tfrac18)$, so
$t=(\tfrac14,\tfrac14,\tfrac16,\tfrac1{12},\tfrac18,\tfrac18)$ and the atoms
are the rational vectors of \S\ref{app:tetra}.  A triple carrying three
distinct directions and missing the direction $d$ has, by \S\ref{app:tetra},
\[
  \SC-I_3\;=\;3I_3-g_d^{\vphantom{\mathsf{T}}}g_d\tp,
  \qquad
  x\tp\bigl(\SC-I_3\bigr)x\;=\;\sum_{i<j}(x_i-x_j)^{2}
\]
after the relabelling that puts $g_d=(1,1,1)$: a sum of squares, singular on
the line $\R g_d$.  Every entry of the certificate is an integer, and its
kernel is one-dimensional at every one of the twelve such triples.  The
equality locus therefore does not destroy the existence of a certificate; it
destroys its slack.

\subsection{The icosahedral design at
\texorpdfstring{$(6,3)$}{(6,3)}}\label{app:icosa}

Put
\begin{equation}\label{eq:icosaconst}
  \rho:=\frac{3}{\sqrt5},\qquad
  \sigma:=\sqrt{\frac{3-\rho}{2}},\qquad
  \mu:=\sqrt{\frac{3+\rho}{2}},
\end{equation}
so that $\sigma^2+\mu^2=3$ and
$\sigma^2\mu^2=\tfrac{9-\rho^2}{4}=\tfrac{9-9/5}{4}=\tfrac95=\rho^2$, whence
$\sigma\mu=\rho$.  The six atoms are the six diagonals of a regular
icosahedron, in the scaling that makes $\ell_c=3$:
\[
  (\sigma,\mu,0),\ (-\sigma,\mu,0),\
  (0,\sigma,\mu),\ (0,-\sigma,\mu),\
  (\mu,0,\sigma),\ (\mu,0,-\sigma),
  \qquad t_c=\tfrac16 .
\]
Each coordinate receives $2(\sigma^2+\mu^2)=6$ from the six outer products and
each off-diagonal entry cancels in pairs, so
$\sum_c t_c g_c^{\vphantom{\mathsf{T}}}g_c\tp=I_3$.  Every distinct pairing is
$\pm\sigma\mu$ or $\mu^2-\sigma^2$, and $\mu^2-\sigma^2=\rho=\sigma\mu$ by
\eqref{eq:icosaconst}; hence
\begin{equation}\label{eq:icosapair}
  \ell_c=3,\qquad \langle g_c,g_d\rangle^2=\tfrac95 \quad(c\ne d).
\end{equation}
For $C=\{0,2,4\}$ all three pairings equal $+\rho$, so
$\SC=(3-\rho)I_3+\rho J_3$, with spectrum $\{3-\rho,3-\rho,3+2\rho\}$ and
\[
  \lmin(S_{\{0,2,4\}})=3-\frac{3}{\sqrt5}=1.6584\ldots>1 :
\]
\begin{sloppypar}
the icosahedral design dominates, and dominates strictly; the margin is
pinned two-sidedly to $2-3/\sqrt5$, which lies in $(\tfrac12,\tfrac23)$
(\texttt{Gtz.icosa\allowbreak Design\_rayleigh\_floor},
\texttt{Gtz.icosa\allowbreak Design\_rayleigh\_attained},
\texttt{Gtz.icosa\_margin\_window}).  By \eqref{eq:icosapair} every triple
carries the same trace and the same pairing modulus, so the verdict depends
only on the sign of the product of the three pairings, and exactly ten of
the twenty triples dominate.  That count is the one assertion of this
appendix with no theorem behind it: the development records it in a
docstring, not a statement.  The design is used in
Appendix~\ref{app:graveyard} as the one on which the two-moment criterion is
silent.
\end{sloppypar}
Figure~\ref{fig:icosa} draws the design and two of its triples.

\begin{figure}[ht]
  \centering
  % The icosahedral design at (6,3), and the spectra of two of its twenty
% triples.  Orthographic projection of R^3 with azimuth 30.5 and elevation 45
% degrees, scaled by 1.3cm:
%   X = 1.3(-0.5075 x_1 + 0.8616 x_2)
%   Y = 1.3(-0.6093 x_1 - 0.3589 x_2 + 0.7071 x_3)
% The twelve vertices are the cyclic permutations of (0, +-1, +-phi) with
% phi = 1.6180340, listed below in six antipodal pairs; the atom g_c is
% sigma = 0.9105930 times the vertex it points at, so every arrowhead stops
% short of its vertex dot.
\begin{tikzpicture}[
  x=1cm, y=1cm,
  vec/.style={-{Latex[length=2mm,width=1.5mm]}, line width=0.75pt, black!55},
  heavy/.style={-{Latex[length=2mm,width=1.5mm]}, line width=1.1pt},
  edge/.style={line width=0.45pt, black!65},
  hidden/.style={line width=0.45pt, black!45, dashed,
                 dash pattern=on 1.4pt off 1.2pt},
  chord/.style={line width=0.55pt, black!55, dash pattern=on 3.6pt off 1.8pt},
  tag/.style={font=\footnotesize, fill=white, inner sep=0.6pt},
  every node/.style={font=\small}]

% ---- the twelve vertices, in six antipodal pairs -----------------------
\coordinate (v5)  at ( 1.1526,-1.5469);   % ( 1, phi,  0)
\coordinate (v8)  at (-1.1526, 1.5469);   % (-1,-phi,  0)
\coordinate (v6)  at ( 2.4722, 0.0372);   % (-1, phi,  0)
\coordinate (v7)  at (-2.4722,-0.0372);   % ( 1,-phi,  0)
\coordinate (v1)  at ( 1.1201, 1.0208);   % ( 0,  1, phi)
\coordinate (v4)  at (-1.1201,-1.0208);   % ( 0, -1,-phi)
\coordinate (v2)  at (-1.1201, 1.9539);   % ( 0, -1, phi)
\coordinate (v3)  at ( 1.1201,-1.9539);   % ( 0,  1,-phi)
\coordinate (v9)  at (-1.0676,-0.3623);   % ( phi,  0,  1)
\coordinate (v12) at ( 1.0676, 0.3623);   % (-phi,  0, -1)
\coordinate (v11) at (-1.0676,-2.2008);   % ( phi,  0, -1)
\coordinate (v10) at ( 1.0676, 2.2008);   % (-phi,  0,  1)
\coordinate (O)  at ( 0,0);

% ---- the six atoms, each sigma times the vertex it points at -----------
\coordinate (g0)  at ( 1.0495,-1.4086);   % g_0 = sigma * ( 1, phi,  0)
\coordinate (g1)  at ( 2.2512, 0.0338);   % g_1 = sigma * (-1, phi,  0)
\coordinate (g2)  at ( 1.0200, 0.9295);   % g_2 = sigma * ( 0,  1, phi)
\coordinate (g3)  at (-1.0200, 1.7792);   % g_3 = sigma * ( 0, -1, phi)
\coordinate (g4)  at (-0.9721,-0.3299);   % g_4 = sigma * ( phi,  0,  1)
\coordinate (g5)  at (-0.9721,-2.0040);   % g_5 = sigma * ( phi,  0, -1)

% ---- 12 hidden edges, the face carrying g_0, g_2, g_4, then the 18
% visible edges ---------------------------------------------------------
\draw[hidden]
  (v2)--(v8) (v3)--(v4) (v3)--(v12) (v4)--(v7) (v4)--(v8)
  (v4)--(v11) (v4)--(v12) (v6)--(v12) (v7)--(v8) (v8)--(v10)
  (v8)--(v12) (v10)--(v12);
\path[fill=black!14] (v5)--(v1)--(v9)--cycle;
\draw[edge]
  (v1)--(v2) (v1)--(v5) (v1)--(v6) (v1)--(v9) (v1)--(v10)
  (v2)--(v7) (v2)--(v9) (v2)--(v10) (v3)--(v5) (v3)--(v6)
  (v3)--(v11) (v5)--(v6) (v5)--(v9) (v5)--(v11) (v6)--(v10)
  (v7)--(v9) (v7)--(v11) (v9)--(v11);

% ---- the three chords joining the tips of g_1, g_3, g_5, which bound no
% face --------------------------------------------------------------------
\draw[chord] (v6)--(v2) (v6)--(v11) (v2)--(v11);

% ---- the six atoms as arrows from the centre ---------------------------
\draw[heavy] (O)--(g0) (O)--(g2) (O)--(g4);
\draw[vec]   (O)--(g1) (O)--(g3) (O)--(g5);
\fill (O) circle (1.1pt);
\fill[black!70]
  (v1) circle (1.0pt) (v2) circle (1.0pt) (v3) circle (1.0pt)
  (v5) circle (1.0pt) (v6) circle (1.0pt) (v7) circle (1.0pt)
  (v9) circle (1.0pt) (v10) circle (1.0pt) (v11) circle (1.0pt);
\fill[black!35]
  (v4) circle (1.0pt) (v8) circle (1.0pt) (v12) circle (1.0pt);
\node[tag, anchor=north west] at ( 1.2721,-1.7073) {$g_0$};
\node[tag, anchor=west] at ( 2.6722, 0.0402) {$g_1$};
\node[tag, anchor=south west] at ( 1.2679, 1.1555) {$g_2$};
\node[tag, anchor=south east] at (-1.2196, 2.1274) {$g_3$};
\node[tag, anchor=east] at (-1.2570,-0.4266) {$g_4$};
\node[tag, anchor=north east] at (-1.1549,-2.3807) {$g_5$};

% ---- the two spectra, on one eigenvalue axis; the eigenvalue lambda is
% drawn at abscissa 0.6 lambda -------------------------------------------
\begin{scope}[shift={(5.55,-0.95)}]
  \draw[line width=0.5pt, black!70] (-0.06,0)--(3.7600,0);
  \foreach \lev in {0,1,2,3,4,5,6}{
    \pgfmathsetmacro{\xx}{0.6*\lev}
    \draw[line width=0.5pt] (\xx,0)--(\xx,-0.07);
    \node[anchor=north, font=\footnotesize, inner sep=2pt]
      at (\xx,-0.07) {\lev};}
  \draw[dashed, dash pattern=on 1.7pt off 1.5pt, line width=0.5pt]
        (0.6000,0)--(0.6000,2.7200);
  \node[anchor=south, font=\footnotesize] at (0.6000,2.7000) {$1$};
  \node[anchor=south, font=\small]
    at (1.8000,3.1200) {$\operatorname{spec}\SC$};
  % C = {0,2,4}: spectrum 1.6584, 1.6584, 5.6833
  \fill[black] (0.9950,2.3000) circle (1.7pt);
  \fill[black] (0.9950,2.4600) circle (1.7pt);
  \fill[black] (3.4100,2.3000) circle (1.7pt);
  \node[anchor=east, align=right, font=\footnotesize, inner sep=0pt]
    at (-0.26,2.3000)
    {$C=\{0,2,4\}$\\[3pt] $\tau=+\rho^{3}$\\[3pt] $\lmin=3-3/\sqrt5$};
  % C = {1,3,5}: spectrum 0.3167, 4.3416, 4.3416
  \fill[black!55] (0.1900,0.6200) circle (1.7pt);
  \fill[black!55] (2.6050,0.6200) circle (1.7pt);
  \fill[black!55] (2.6050,0.4600) circle (1.7pt);
  \node[anchor=east, align=right, font=\footnotesize, inner sep=0pt]
    at (-0.26,0.6200)
    {$C=\{1,3,5\}$\\[3pt] $\tau=-\rho^{3}$\\[3pt] $\lmin=3-6/\sqrt5$};
  % what the two rows share
  \node[anchor=north west, align=left, font=\footnotesize, inner sep=0pt]
    at (-2.16,-0.52)
    {$\rho=3/\sqrt5$, \quad $\tau=\langle g_a,g_b\rangle
     \langle g_a,g_c\rangle\langle g_b,g_c\rangle$\\[3.5pt]
     $\tr\Gram_C=9$ and $\langle g_c,g_d\rangle^{2}=\tfrac95$ at
     all twenty triples\\[3.5pt]
     $\gamma(C)=8\tau-\tfrac{104}{5}<0$ at all twenty triples};
\end{scope}
\end{tikzpicture}

  \caption{The icosahedral design at $(6,3)$, in the constants of
  \eqref{eq:icosaconst}: six atoms along the six diagonals of a regular
  icosahedron, each drawn as $\sigma$ times the vertex it points at.  All
  twenty triples carry the trace $9$ and the single pairing modulus
  $\rho$, and ten of them dominate.  The verdict is the sign of the
  product $\tau$ of the three pairings.  The heavy atoms point at the
  vertices of the shaded face, the light ones at three vertices joined by
  no edge.  By Propositions~\ref{prop:icosa} and~\ref{prop:momentsharp} no
  criterion in the trace and the determinant certifies the triple that
  dominates.}
  \label{fig:icosa}
\end{figure}

\subsection{Active-family census at
\texorpdfstring{$(6,3)$ and $(7,3)$}{(6,3) and (7,3)}}\label{app:census}

By \eqref{eq:coverage} the family $A^{+}$ of active $k$-subsets
carrying a positive multiplier covers the index set.  The covering families
of $3$-subsets are counted here up to the symmetric group of the index set
and listed by cardinality.  At $m=6$ a covering family has at
least two members and at most twenty; Table~\ref{tab:census} lists the
counts, whose total is $2102$.  The single class of cardinality two is the
partition branch that Corollary~\ref{cor:e3} leaves undecided; setting
it aside leaves $2101$.  By Proposition~\ref{prop:carath} the support of the
multiplier vector may be taken of size at most $15$ here, so only the columns
up to $15$ bear on
Conjecture~\ref{conj:sixthree}; they account for $2069$ of the $2102$ classes,
and for $2068$ of the $2101$ with at least three members.

At $m=7$ a covering family has at least three members, and the counts are
\[
\begin{array}{c|cccccccc}
\lvert A^{+}\rvert & 3&4&5&6&7&8&9&10\\ \hline
\text{classes} & 3&17&94&415&1599&5308&15397&39081\\[4pt]
\lvert A^{+}\rvert & 11&12&13&14&15&16&17&18\\ \hline
\text{classes} & 87347&172684&303116&473733&660907&824389&920439&920443\\[4pt]
\lvert A^{+}\rvert & 19&20&21&22&23&24&25&26\\ \hline
\text{classes} & 824409&660949&473827&303277&172933&87659&39433&15709\\[4pt]
\lvert A^{+}\rvert & 27&28&29&30&31&32&33&34\\ \hline
\text{classes} & 5557&1760&509&137&38&10&3&1
\end{array}
\]
together with a single class of cardinality $35$,
with total $7\,011\,184$.  By Proposition~\ref{prop:carath} the bound at
$(7,3)$ is $19$, and the classes of cardinality at most
$19$ already number $5\,249\,381$.  Enumeration at this size is not a route;
Problem~\ref{prob:seventhree} asks instead for an argument uniform in
the active family.
